\documentclass[11pt]{amsart}
\usepackage[localtheorems]{raeez-math-template}

\providecommand{\C}{\mathbb{C}}
\providecommand{\bC}{\mathbb{C}}
\providecommand{\R}{\mathbb{R}}
\providecommand{\Z}{\mathbb{Z}}
\providecommand{\dbar}{\bar{\partial}}
\providecommand{\g}{\mathfrak{g}}
\providecommand{\Obs}{\mathrm{Obs}}
\providecommand{\Sym}{\mathrm{Sym}}
\providecommand{\CE}{\mathrm{CE}}
\providecommand{\chir}{\mathrm{chir}}
\providecommand{\Ran}{\mathrm{Ran}}
\providecommand{\Hom}{\mathrm{Hom}}
\providecommand{\cL}{\mathcal{L}}
\providecommand{\cM}{\mathcal{M}}
\providecommand{\cE}{\mathcal{E}}
\providecommand{\cD}{\mathcal{D}}
\providecommand{\cO}{\mathcal{O}}
\providecommand{\kappach}{\kappa_{\mathrm{ch}}}

\theoremstyle{plain}
\newtheorem{theorem}{Theorem}[section]
\newtheorem{proposition}[theorem]{Proposition}
\newtheorem{lemma}[theorem]{Lemma}
\theoremstyle{definition}
\newtheorem{definition}[theorem]{Definition}
\newtheorem{construction}[theorem]{Construction}
\theoremstyle{remark}
\newtheorem{remark}[theorem]{Remark}

\title{Chiral Chevalley--Eilenberg cochains for 6D hCS}
\author{Wave 13 E4/20}
\date{}

\begin{document}
\maketitle

\section{The target}
For a dg Lie algebra $\g$ on a complex $3$-fold $X$, we construct a
cochain complex $\CE^*_{\dbar,\chir}(\cL, \cM)$ computing derived
deformations of $\cL$-modules in the chiral--Dolbeault bicomplex, and
identify it with $\Obs_{\mathrm{hCS}}$ for the holomorphic Lie algebroid
$\cE = \Omega^{0,\bullet}(\C^3, \g)$.

\section{Cycle 1: Classical CE}
For a Lie algebra $\g$ over $k$ with module $M$, classical CE cochains
are
\[
 \CE^n(\g, M) = \Hom_k\!\bigl(\Lambda^n \g, M\bigr), \qquad
 d_{\CE} f(x_1, \ldots, x_{n+1})
  = \sum_{i<j} (-1)^{i+j} f\!\bigl([x_i, x_j], \ldots \widehat{x_i}
     \ldots \widehat{x_j} \ldots\bigr)
  + \sum_i (-1)^{i+1} x_i \cdot f(\ldots \widehat{x_i} \ldots).
\]
\textbf{Attack.} The extension to chiral requires replacing
$\Lambda^n\g \subset \g^{\otimes n}$ by a structure supported on the
Ran space, and the antisymmetrisation by a $\Sigma_n$-equivariant
pushforward. \textbf{Heal.} This is exactly Beilinson--Drinfeld's
chiral symmetric power $\Sym^n_{\chir}\cL$ on $\Ran X$
(\cite[\S3.4]{BD04}).

\section{Cycle 2: Chiral CE}
Let $\cL$ be a chiral Lie algebra on a smooth curve (or product of
curves) $X$: a right $\cD_X$-module with chiral bracket
$\mu: j_*j^*(\cL \boxtimes \cL) \to \Delta_*\cL$ satisfying
skew-symmetry and chiral Jacobi.
\begin{definition}[Chiral CE cochains]
\label{def:chiral-ce-cochains}
Fix a chiral $\cL$-module $\cM$. The chiral CE cochain complex is
\[
 \CE^n_{\chir}(\cL, \cM)
   := \Hom_{\cD_{X^n}}\!\bigl(\Sym^n_{\chir}\cL,\; \Delta^{(n)}_*\cM\bigr),
\]
where $\Sym^n_{\chir}\cL$ is the BD chiral symmetric power supported on
pairwise-distinct configurations, and $\Delta^{(n)}: X \hookrightarrow
X^n$ is the diagonal. The chiral differential
$d_{\chir}: \CE^n_{\chir} \to \CE^{n+1}_{\chir}$ is
\[
 (d_{\chir} f)(l_0, \ldots, l_n)
  = \sum_{i<j} (-1)^{i+j}\, f\!\bigl(\mu_{ij}(l_i, l_j),\,
       \ldots \widehat{l_i} \ldots \widehat{l_j} \ldots\bigr)
  \;+\; \sum_i (-1)^{i+1}\, \rho_i \cdot f(\ldots \widehat{l_i} \ldots),
\]
with $\mu_{ij}$ the chiral bracket on positions $(i,j)$ pushed to the
diagonal and $\rho_i$ the chiral action on $\cM$. The Jacobi identity
on $\Ran^{\leq 3}$ yields $d_{\chir}^2 = 0$.
\end{definition}
\textbf{Attack.} Does the $\cD$-module Hom detect sufficiently refined
information? \textbf{Heal.} Yes: by
\cite[Prop.~3.4.11]{BD04}, $\Sym^n_{\chir}\cL$ is a projective
generator for chiral $n$-ary operations, and
$\CE^*_{\chir}(\cL,\cL) \simeq \Hom_{\chir}(\cL[-1], \cL[-1])$ in the
derived category of chiral Lie algebras. This recovers
Beilinson--Drinfeld's chiral-envelope bar duality
$\CE_*^{\chir}(\cL) \simeq B(U^{\chir}(\cL))$ upon dualising (compare
\texttt{cy\_to\_chiral.tex}, \texttt{prop:bar-ce-chiral}).

\section{Cycle 3: Holomorphic variant}
On a complex $3$-fold $X$, replace $\cD_X$-modules by $\cD_X$-modules
in Dolbeault degrees: take $\cL \in \Omega^{0,\bullet}(X, \g) \otimes
\Omega^{\bullet}_{\mathrm{hol}}$, i.e.\ a graded sheaf carrying both a
$\dbar$-differential and a chiral bracket along the holomorphic
directions.
\begin{definition}[Dolbeault--chiral CE]
\label{def:dolbeault-chiral-ce}
\[
 \CE^*_{\dbar, \chir}(\cL, \cM)
  := \Bigl(\,\bigoplus_{n \geq 0}\;
     \Omega^{0, \bullet}\!\bigl(X^n,\; \Hom(\Sym^n_{\chir}\cL, \cM)\bigr)^{\Sigma_n}
     ,\; \dbar + d_{\chir}\Bigr),
\]
a double complex with total differential $D = \dbar + d_{\chir}$.
Totalisation yields a single complex; $D^2 = 0$ because $\dbar^2 = 0$,
$d_{\chir}^2 = 0$, and $\dbar$ commutes with the chiral bracket (which
is purely holomorphic).
\end{definition}
\textbf{Attack.} Compatibility $[\dbar, d_{\chir}] = 0$ is
nontrivial when $\cL$ has nonholomorphic components of its bracket.
\textbf{Heal.} For $\cL = \Omega^{0,\bullet}(X,\g)$ with pointwise
$\g$-bracket on coefficients, the chiral bracket is the
$\delta$-function bracket $\mu(\alpha \otimes \beta) =
[\alpha,\beta]_\g \cdot \delta_\Delta$; both operations commute with
$\dbar$ since the bracket acts on coefficients and $\dbar$ on
differential form type.

\section{Cycle 4: 6D hCS application}
For holomorphic Chern--Simons on $\C^3$ with structure Lie algebra
$\g$, the Costello BV theory (\cite{CostelloGaiotto}) has fields
$\cA \in \Omega^{0,\bullet}(\C^3, \g)[1]$ and action
$S = \int_{\C^3} \Omega \wedge \langle \cA, \dbar \cA + \tfrac{1}{3}
[\cA,\cA] \rangle$ with $\Omega$ the holomorphic volume form.
\begin{proposition}[hCS observables as chiral CE]
\label{prop:hcs-obs-chiral-ce}
Let $\cE := \Omega^{0,\bullet}(\C^3, \g)$ with $\dbar$-differential and
pointwise bracket, viewed as a holomorphic Lie algebroid. Then
\[
 \Obs_{\mathrm{hCS}}^{cl}(\C^3)
   \;\simeq\; \CE^*_{\dbar, \chir}\!\bigl(\cE,\; \cO_{\C^3}\bigr),
\]
and quantum observables are the BV-deformation computed by adding the
Costello counterterm $\int \Omega \wedge \mathrm{tr}(\mathrm{ad}_\cA)$
to $d_{\chir}$ in the leading $\hbar$-bracket.
\end{proposition}
\textbf{Attack.} Is $\cE$ really chiral in the BD sense on $\C^3$?
\textbf{Heal.} $\C^3$ is not a curve, but BD's construction extends to
higher-dimensional Ran spaces via Francis--Gaitsgory
\cite[\S2]{FrancisGaitsgory} and the BD formalism on $X^n$ for $\dim
X > 1$ yields an $E_3$-chiral structure (\texttt{CLAUDE.md}: at
$d \geq 3$, $A$ is $E_1$; $E_2/E_3$ live elsewhere). Here the chiral
bracket is $E_1$ along the holomorphic direction of each factor of
$\C^3$, and the $\dbar$-enhanced complex captures the full $E_3$-chiral
observable algebra. This matches Costello's original identification
$\Obs \simeq \CE^*_{\dbar}(\cE)$ (E2 of this wave) and refines it by
making the chiral structure explicit.

\section{Cycle 5: Comparison with Francis 2013}
Francis \cite[\S3]{Francis13} defines $E_n$-factorisation homology
$\int_M A$ for an $E_n$-algebra $A$ and framed $n$-manifold $M$.
\begin{proposition}[Koszul duality]
\label{prop:koszul-factorisation}
For $\cE = \Omega^{0,\bullet}(\C^3,\g)$,
\[
 \int_{\C^3} U^{\chir}(\cE)
   \;\simeq\; \CE^*_{\dbar, \chir}(\cE, \C)^\vee
\]
as $E_3$-algebras, with $U^{\chir}$ the chiral envelope and $(-)^\vee$
linear dual. This is the Francis--Ayala $E_n$-Koszul duality between
$E_n$-enveloping algebras and $E_n$-CE cochains specialised to the
chiral context of Kapranov--Vasserot \cite[\S5]{KapranovVasserot14}.
\end{proposition}
\textbf{Attack.} Francis' theorem requires $A$ augmented and
connective. \textbf{Heal.} $U^{\chir}(\cE)$ is augmented via
$\cE \to 0$ and connective in the Dolbeault grading; the pair
$(U^{\chir}\cE, \CE^*_{\dbar,\chir}(\cE,\C))$ are Koszul duals
(Francis--Gaitsgory proof, extended to $d = 3$ via the
Kapranov--Vasserot chiral Hall category).

\section{Consequence: $\kappa_{\mathrm{ch}}$ from CE cohomology}
For compact CY$_3$ $X$ with $\cE = \Omega^{0,\bullet}(X,\g)$,
\[
 \kappa_{\mathrm{ch}}(A_X) = \chi\!\bigl(\CE^*_{\dbar,\chir}(\cE, \cO_X)\bigr)
   = \sum_q (-1)^q h^{0,q}(X) \cdot \dim \CE^*(\g)_{\text{hor}},
\]
recovering the Hodge-supertrace identification on compact CY$_d$
(CLAUDE.md essential constants). For $\g$ abelian and $X = \C^3$ the
ghost-number weighted character matches
\texttt{wave12\_a5\_hCS\_BV\_BRST\_explicit.tex}.

\begin{thebibliography}{9}
\bibitem{BD04} A.~Beilinson and V.~Drinfeld, \emph{Chiral Algebras},
  AMS Colloquium Publications 51, 2004, \S3.4--3.7.
\bibitem{Francis13} J.~Francis, \emph{The tangent complex and Hochschild
  cohomology of $E_n$-rings}, Compositio Math.\ 149 (2013), 430--480,
  \S3.
\bibitem{KapranovVasserot14} M.~Kapranov and E.~Vasserot,
  \emph{The cohomological Hall algebra of a surface and factorization
  cohomology}, arXiv:1901.07641 / Kapranov--Vasserot notes 2014, \S5.
\bibitem{FrancisGaitsgory} J.~Francis and D.~Gaitsgory,
  \emph{Chiral Koszul duality}, Selecta Math.\ 18 (2012), 27--87.
\bibitem{CostelloGaiotto} K.~Costello and D.~Gaiotto, \emph{Twisted
  supergravity and its quantization}, arXiv:1812.01110.
\end{thebibliography}

\end{document}
